\documentclass[12pt]{article}
\usepackage[utf8]{inputenc}
\usepackage[T1]{fontenc}
\usepackage{hyperref}
\title{Tumtum Language Resources for K-12 Curriculum Development}
\author{Compiled by Claude Sonnet 5 (claude-sonnet-5) for Andr\'as Kornai}
\date{}
\begin{document}
\maketitle

\section{Introduction}
Tumtum (ISO 639-3 tbr, Glottocode tumt1243) is a Kadu (Kadugli-Krongo) language of South Kordofan, in the Nuba Mountains of Sudan. Its dialects are Karondi, Talassa and Tumtum proper. Speaker counts differ: 17,000 (Wikipedia, 2022) and 24,000 (Ethnologue 2016, via Joshua Project). UNESCO classifies it as severely endangered, according to Wikipedia. The Nilo-Saharan placement of the Kadu family is contested.

\section{Primary Educational Resources}
No primers, textbooks or school materials were found in search results.

\section{Linguistic and Reference Materials}
\begin{itemize}
\item Wikipedia, Tumtum language: \url{https://en.wikipedia.org/wiki/Tumtum_language}
\item Wikipedia, Kadu languages: \url{https://en.wikipedia.org/wiki/Kadu_languages}
\item Leiden University repository items (titles not shown in results; possibly Kadu-related scholarship, not verified): \url{https://scholarlypublications.universiteitleiden.nl/access/item%3A2729430/view} and \url{https://scholarlypublications.universiteitleiden.nl/access/item%3A2728068/view}
\end{itemize}

\section{Religious and Community Resources}
\begin{itemize}
\item Joshua Project, Tumtum in Sudan (lists Bible translation status as ``Translation Needed''): \url{https://joshuaproject.net/people_groups/15638}
\end{itemize}

\section{Digital and Audio Resources}
No audio, video or corpus resources were found in search results.

\section{Curriculum Development Context}
No evidence of schooling in Tumtum or of a standard orthography was found. Conflict in the Nuba Mountains limits access to the area. A curriculum would need new documentation work with the community.
\end{document}